Thus $\delta=1-g$ and $m=\lambda+1\ge2$.  If $y=0$, the complete-word row of
Lemma~\ref{lem:alt-lift} gives a source-zero constant-tail state instead.
\end{proof}

\begin{lemma}[LOSS-anchored lift and factor module]
\label{lem:loss-anchored-lift}
Let $y$ be an already retained finite LOSS token.  Suppose a continuing DRAW
obligation or adjacent factor frame is attached, by exact coordinates, to a
constant-tail state
\[
 x=Q_m^d(3^kJ(y)),\qquad m\ge1,\quad k\ge0.
\]
Then the attached normalization has only two kinds of outcome-compatible
macro-exit:
\begin{enumerate}[label=(\alph*)]
\item $y$ is replaced by one of its actual WIN children, or by a certified
      descendant of such a child, and the proof-token rank is strict;
\item $y$ is retained and an exact tail/factor counter decreases.  No
      unrelated finite endpoint is inserted into a ranked token multiset.
\end{enumerate}
In particular the module never exits as a fresh free obligation at an
unrelated numerical cursor.
\end{lemma}

\begin{proof}
Both ordinary children
\[
 c_A=A(y),\qquad c_B=B(y)
\]
of the LOSS state $y$ are WIN and satisfy
$h(c_A),h(c_B)\le h(y)-1$.  While a constant tail has length at least three,
Lemma~\ref{lem:long-tail} removes one or two tail bits and merely multiplies
the odd coefficient by $3$; its underlying three-free source remains $y$.
The valuation-two transfer and the canonical $A$ continuation therefore stay
inside the same attached module while the exact tail counter decreases.  No
finite endpoint encountered during this countdown is inserted into a ranked
multiset merely because its outcome is known.

It remains to audit a factor boundary.  The universal current/next-level
coupling reduces its first factor-free signed source over the finite source
$y$ to
\[
 9J(y)+1-2f=2^vJ(t).
\]
The exact source trichotomy is
\[
 t=\begin{cases}
   Q_r^\epsilon(J(c_B))\quad(r\ge1),&v=1,\\
   A(c_A),&v=2,\\
   B(c_A),&v\ge3.
  \end{cases}
\]
For $v=1$ this is Lemma~\ref{lem:alt-lift} applied to $A(y)$; for
$v\ge2$ it is the ordinary source-boundary transition at $A(y)$.  Thus the
$v=1$ row is already an exact lift over the proper WIN child $c_B$, whereas
$v\ge2$ is an actual child of the proper WIN child $c_A$.  Before any new
normalizer is installed we may therefore make the strict retained-token
replacement $y\mapsto c_B$ in the first row or $y\mapsto c_A$ in the other
two rows.  The returned cursor itself is not inserted as a token.

The consecutive-factor identity is equally important: when the current
valuation is at least two, the next factor source is an exact lift over the
already displayed $t$; when it is one, the next factor source is exactly one
ordinary selected-source step from that first lift.  Hence a DRAW choosing a
raw rather than signed exit cannot bypass the strict child token.  If powers
of three remain at tail exponent at least two,
Lemma~\ref{lem:fixed-tail-factor} removes them one level at a time or exposes
one of its typed finite side diamonds.  At exponent one the same two-level
factor coupling is used instead; the reverse-factor lemma is never invoked
there.  Consequently equal-token routing has a decreasing exact natural
counter, and every genuine restart is preceded by the displayed child-token
replacement.  This is the LOSS-source provenance lemma of the older local
proof, rewritten here so that it is not hidden inside the global assembly.
\end{proof}

\begin{lemma}[The $D=2$ exact lift remains LOSS-anchored]
\label{lem:D2-pure}
The $D=2$ row of Lemma~\ref{lem:short-tail} is an instance of
Lemma~\ref{lem:loss-anchored-lift} with retained LOSS token $y$.  Therefore
it either makes a strict descendant-token transition before control restarts,
or stays in a finite exact tail/factor countdown over the same $y$.  It never
uses a new seed and never becomes a free obligation at the returned cursor
$w$.
\end{lemma}

\begin{proof}
For $y=0$ this is Lemma~\ref{lem:zero-source}.  For $y>0$ the preceding
calculation gives $w=Q_m^{1-g}(J(y))$ with $m\ge2$.  If the first factor
numerator is $N_1=2J(w)$, then consecutive factor levels satisfy
\[
 N_2=3N_1-2(1-2g)
    =2\bigl(3J(w)+1-2(1-g)\bigr).
\]
Thus the next signed source is exactly the ordinary source-boundary
transition of this lift.  Lemma~\ref{lem:loss-anchored-lift} applies from
that point and keeps $y$ as the retained provenance until a strict
child-token event occurs.  A later marked-tail return is therefore either
after that strict event or still inside the same decreasing attached
countdown; there is no unranked same-fibre reseed.
\end{proof}

\begin{lemma}[Post-payment shortest marked lift arithmetic]
\label{lem:attached-lift}
Consider the $D=3$ row of Lemma~\ref{lem:marked-tail-pay} after the strict
replacement $y\mapsto r$ has already been made. Thus
\[
 b=Q_3^g(C),\quad q=A(b)\text{ WIN},\quad y=B(b)\text{ LOSS},\quad
 r=B(q)=B(y)\text{ WIN}.
\]
Since $q$ and $r$ are WIN, its other child $\ell=A(q)$ is LOSS. The two
children of $\ell$ are WIN and are proper descendants of $q$.

The first raw exit of the marked factor scan is an exact lift
\[
 S=Q_m^g(J(r)),\qquad m\ge2.
\]
For an arbitrary A-selecting lift cursor $x=Q_m^g(a)$, with
$e=1-g=\alpha(x)$, the common side $d$ of $O(x,e)$ satisfies
\[
 d=\begin{cases}
 9a-g,&m=2,\\
 R(Q_1^g(9a)),&m=3,\\
 Q_{m-3}^g(9a),&m\ge4,
 \end{cases}
\]
and the source transition at $A(x)$ has respectively valuation $1$, at least
$3$, and $2$. Hence every same-provenance canonical recycle with $m\ge4$
has the exact update
\[
 (m,a)\longmapsto(m-3,9a),
\]
which preserves the underlying three-free source and lowers the natural lift
counter. Every noncanonical exit remains an exact raw/signed boundary over
$r$ or over an actual child of the transported lower pair. No temporary
finite endpoint is promoted to the retained token multiset.
\end{lemma}

\begin{proof}
The identity $S=Q_m^g(J(r))$, $m\ge2$, is the exact appended-suffix
calculation at the $D=3$ boundary. The three-row formula follows by applying
Lemma~\ref{lem:long-tail} to $A(x)$ and by using the shared exponent-two/one
boundary at $m=2,3$. For $m\ge4$ the tail has at least four constant bits,
so one canonical A/B$_2$ block removes three of them and multiplies the odd
coefficient by $9$, giving the displayed update.

The retained geometry is independent of the value of $m$. The WIN state $q$
has the WIN child $r$, hence $\ell=A(q)$ is LOSS; both children of $\ell$
are WIN proper descendants of $q$. Repeated common-child steps transport this
pair componentwise while the raw exponent is reduced by three. Therefore the
factor/signed exits stay on the exact boundary generated by $r$ and its actual
children. This lemma asserts provenance transport only.  The proof-height
payment at the terminal classes is proved separately after this arithmetic
transport has been established.
\end{proof}

\begin{lemma}[Attached $B$-selecting exponent-one lifts descend]
\label{lem:attached-b}
Let $r>0$, $d\in\{0,1\}$, and put
\[
 S=Q_1^d(J(r)),\qquad e=1-d.
\]
Then $e$ is the $B$-selecting phase at $S$.  If $k=B(S)>0$, its state source
satisfies
\[
 \boxed{\rho(k)<r}.
\]
Hence a $B$-selecting obligation attached to this exact exponent-one lift
transfers, by Lemma~\ref{lem:b-diamond}, either to the terminal source or to
an adjacent frame with a strictly smaller retained nested source.  It never
becomes a fresh free obligation at the numerical cursor $S$.
\end{lemma}

\begin{proof}
A one-bit tail of bit $d$ has A-selecting phase $d$, so $e=1-d$ is
B-selecting.  The universal contraction and source bounds give
\[
 k=B(S)\le\left\lfloor\frac{3S+1}{4}\right\rfloor,
 \qquad
 \rho(k)\le\frac{k-1}{6}\le\frac{S-1}{8}.
\]
Since $J(r)\le3r+2$, one has $S\le6r+4$, and therefore
$(S-1)/8<r$ for $r\ge2$.  For $r=1$ the two lifts are $S=10,9$;
their B-children are $7,3$, both of state source $0$.  Thus the inequality
holds in every positive case.  Lemma~\ref{lem:b-diamond} supplies the exact
transferred adjacent frame over $k$.
\end{proof}

\begin{lemma}[Every terminal macro pays a retained token]
\label{lem:terminal-pay}
Suppose a terminal row is reached by the post-payment $D=3$ lift transport of
Lemma~\ref{lem:attached-lift}.  Let $\ell=A(q)$ be the retained LOSS token
whose two WIN children form the transported lower pair; in the zero-recycle
endpoint the already retained marked LOSS token is $y$.  The common-child transport is componentwise: after every
three-bit recurrence both transported WIN states are proper descendants of
the corresponding incoming WIN states.  At the terminal barrier they are
the complete child set of one reconstructed LOSS state $Z_n$.

Every terminal exponent-one, exponent-two, or exponent-three restart is
therefore preceded by a strict replacement of retained provenance.  More
precisely:
\begin{enumerate}[label=(\alph*)]
\item the zero-recycle two-bit row $m=2$ either exposes a boundary WIN at
      least two proof levels below its boundary-WIN source, or enters an exact
      adjacent DRAW frame over that source's actual LOSS child one level
      lower;
\item the exponent-one row first reconstructs a canonical LOSS token $L$
      with $h(L)\le h(\ell)-2$ and only then enters its $B$-selecting gate;
\item every positive-recycle exponent-two or exponent-three row reconstructs
      a LOSS token $Z_n$ with $h(Z_n)<h(\ell)$ before the new inner task is
      installed;
\item in the sole zero-recycle exponent-three row, the existing marked LOSS
      token $y$ is replaced by its actual WIN child $r$ with
      $h(r)\le h(y)-1$.
\end{enumerate}
Consequently no terminal row is a pure control reset.
\end{lemma}

\begin{proof}
First isolate the zero-recycle two-bit row, which must not be hidden inside
the positive-recycle inverse-parent formulas.  Let
\[
 x=Q_2^g(a)\text{ WIN},\qquad e=1-g,\qquad O(x,e),
\]
and put
\[
 y=A(x),\quad c=B(x),\quad d=A(y),\quad p=B(y),
\]
where $d$ is the common side of the obligation.  If $d$ is DRAW, then $y$ is non-LOSS, so the WIN state $x$ has the
actual LOSS child $c$ and $h(c)=h(x)-1$.  If $x$ is ordinary, the ordinary
side identity makes $p=B(A(x))$ a child of the LOSS state $c$.  Hence $p$ is
WIN and $h(p)\le h(x)-2$.  The two children of $y$ are precisely $d$ and
$p$; since $d$ is DRAW and $p$ is WIN, outcome recursion now gives that $y$
is DRAW.  Thus $y\to p$ is the required lower boundary.  If $x$ is
exceptional, Lemma~\ref{lem:exceptional-side} gives the exact adjacent DRAW
frame over the actual LOSS child $c$.

If $d$ is WIN, apply the finite-source export lemma
(Lemma~\ref{lem:finite-front}) directly to the finite WIN source $x$ and the
A-selecting obligation $O(x,e)$.  It gives a certified proper descendant
before the continuation can restart, except in the unique exceptional
orientation in which the canonical LOSS child is $c=B(x)$; there the exact
exceptional child-pair identity retains an adjacent frame over $c$, with
$h(c)=h(x)-1$.  Thus the $m=2$, zero-recycle row always pays a genuine
descendant token before any restart.

The remaining componentwise claim is the exact inverse-parent calculation for
the shortest marked pair after at least one three-bit transport step (or for
the zero-recycle exponent-three row).  If the first raw lift exponent is $M=3h+n$ with
$n\in\{1,2,3\}$, one common-child recurrence replaces both coordinates by
actual children and lowers their proof heights componentwise; after $h$
recurrences the pair is one of the three terminal barriers.  In the
exponent-two barrier the two WIN states are exactly the children of
\[
 Z_2=Q_2^e(3^{2h-1}J(r))\quad(h\ge1),
\]
and in the positive exponent-three barrier they are the two children of
\[
 Z_3=Q_3^e(3^{2h-1}J(r))\quad(h\ge1).
\]
The LOSS height formula (one plus the maximum of its two child heights)
preserves the componentwise decrement.  In the exponent-one barrier, for
$h\ge2$ the two WIN states are exactly the children of
\[
 Z_1=Q_4^e(3^{2h-3}J(r)),
\]
while for $h=1$ this parent is the original $\ell$.  Hence
\[
 h(Z_1)\le h(\ell)-2h+2,
\]
and the analogous positive-recycle $Z_2,Z_3$ bounds are strict.  These are
exact state identities; no new finite outcome is inferred from a numerical
source label.

For the exponent-one row retain the terminal notation
\[
 a=3^{2h}J(r),\qquad h\ge1,\qquad X=Q_1^g(a),\qquad e=1-g.
\]
The reverse terminal barrier reconstructs a LOSS state $Z_1$.  Its upper
child is
\[
 U=Q_3^e(a/9)=A(Z_1)\quad\text{WIN}.
\]
The two children of $U$ are exactly
\[
 Q_2^e(a/3),\qquad Q_1^e(a/3).
\]
Choose a canonical LOSS child $L$ of this WIN state.  Then
\[
 h(L)=h(U)-1\le h(Z_1)-2\le h(\ell)-2,
\]
(the sharper bound in the positive-recycle rows is not needed here).  With
$c=a/3$ the new terminal source is $X=Q_1^g(3c)$ and its phase $e$ is
$B$-selecting.  Thus the old token has already been replaced strictly before
this gate is normalized.

For completeness, that gate has no hidden unbounded restart.  Put
\[
 W=B(Q_1^e(c))=B(Q_2^e(c)).
\]
Since $L$ is one of the displayed exponent-one/two states, $W$ is an actual
WIN child of $L$ and $h(W)\le h(L)-1$.  If
$\lambda=\operatorname{asl}(9c-g)$ is at least four, its transferred source
$R(9c-g)$ is strictly below $c$.  For $\lambda=1,2,3$ the exact short rows
are respectively: a factor-free lift over $R(W)<c$; a factor-free lift over
the lower WIN token $W$; or the ordinary child $A(W)$.  Hence all three
short rows occur only after the strict source/token payment already displayed.
They may therefore initialize later inner routing data, but none is a
same-rank reseed of a free obligation.

For positive-recycle exponent three, the reconstructed parent can be written
\[
 Z_3=Q_3^e(3^{2h-1}J(r))\quad(h\ge1),
\]
whose two terminal barrier children are WIN; the reverse construction gives
$h(Z_3)\le h(\ell)-2h<h(\ell)$.  For exponent two ($h\ge1$),
\[
 Z_2=Q_2^e(3^{2h-1}J(r))
\]
has the two reconstructed WIN barrier children and likewise
$h(Z_2)\le h(\ell)-2h<h(\ell)$.  These are strict replacements before the
new obligation/factor state is installed.

Finally, in the zero-recycle exponent-three row the marked configuration is
precisely the $D=3$ case of Lemma~\ref{lem:marked-tail-pay}; hence
$r=B(q)=B(y)$ is an actual WIN child of the already present LOSS token $y$.
Thus $h(r)\le h(y)-1$.  Every terminal row therefore pays a source/token
component before any reset.
\end{proof}

\begin{corollary}[The post-payment $D=3$ lift module is well-founded]
\label{cor:short-lift-closed}
After the strict marked-tail token edge at $D=3$, every outcome-compatible
shortest-row lift continuation either makes another strict retained
source/token transition or, with retained provenance fixed, decreases the
exact natural lift exponent until it reaches $m\in\{1,2,3\}$, where
Lemma~\ref{lem:terminal-pay} makes a strict retained-token transition before
any restart. Hence no equal-retained-data cycle is contained in this module.
\end{corollary}

\begin{proof}
Lemma~\ref{lem:attached-lift} gives
$(m,a)\mapsto(m-3,9a)$ for every same-provenance canonical row with
$m\ge4$ and transports the exact lower pair through all temporary factor
exits. Thus $m$ reaches $1,2$ or $3$ in finitely many recurrences unless a
strict source/token exit occurs earlier. At those barriers
Lemma~\ref{lem:terminal-pay} supplies the strict payment. The arithmetic
transport is proved before the terminal payment, so the dependency is
one-way rather than circular.
\end{proof}

\begin{proposition}[Semantic exhaustiveness of the marked router]
\label{prop:P2-closed}
Every actual outcome-compatible DRAW continuation entering the typed routing
system belongs to one of the semantic rows below.  In each row all finite
states inserted into $M$ or $N$ are states already present in the displayed
parent/child geometry; no unrelated finite endpoint is promoted merely because
its outcome is known.

\begin{enumerate}[label=\textnormal{(S\arabic*)},leftmargin=*]
\item \textbf{Boundary or loss-sibling countdown.}  A factor-free frame
\[
 \{Q_{m+1}^e(J(r)),Q_m^e(J(r))\},\qquad m\ge1,
\]
known to contain a DRAW uses Lemma~\ref{lem:long-tail} while $m\ge3$.
The exact tail exponent decreases.  At $m=1,2$ the row is respectively an
obligation/factor boundary or the typed three/two constructor of
Lemma~\ref{lem:hidden-parent}.  A loss-sibling entry of
Definition~\ref{def:entry-controls} carries, in addition, its exact finite
WIN/LOSS witness and the finite local counter $\tau$ supplied by the
constructing diamond.  The fixed-tail and exponent-one predecessor rows
remove one factor/tail level at each reverse normalization step; the
ordinary, exceptional and hidden-parent side rows are one-step exact
attachments.  Thus before leaving (S1) the witness is merely retained (or is
replaced by a proved descendant) while $\tau$ strictly decreases to a
boundary, free obligation, or attached factor/lift row.  No fresh token is
needed inside this countdown.

\item \textbf{Free/outer obligation.}  A free obligation before the first
finite token may use the outer seed only according to
Lemma~\ref{lem:seed-policy}; after $\eta=0$ a first inner witness may use
$\zeta$ once.  Every canonical continuation and valuation-two reset is
exactly one of Proposition~\ref{prop:P1-closed}'s source/seed/descendant/
attached exits.  A reverse edge to a new free obligation is forbidden unless
that proposition has already paid one of those components.

\item \textbf{Finite-source A obligation.}  If its ordinary source is an
outer or inner finite token, Lemma~\ref{lem:finite-front} exports a certified
proper descendant before a later high return.  In the exceptional orientation
the two children are the exact adjacent lifts over the canonical lower LOSS
token given by Lemma~\ref{lem:exceptional-side}.

\item \textbf{Lower factor fork.}  Its lower exponent is exactly $r=1$ or
$r\ge2$.  Lemma~\ref{lem:factor-fork-pay} gives, respectively, a lower
boundary/LOSS token or the source/WIN/LOSS trichotomy.  The only nonterminal
continuation is the typed high-return scan carrying the same finite
provenance.

\item \textbf{Short high return.}  The returned valuation is $v=5$ or
$v=6$, and the constructor carries the finite return token
$c=Q_{v-3}^g(3J(t))$ with its proved outcome.  Lemma~\ref{lem:high-return-pay}(a)
